%% kapa: exam=1 round=22 year=2011 period=2 form=A subject=accounting subject_ko=회계학 seq=17 number=17
%% source: Q-Net 감정평가사 기출문제 (cst00309) artlSeq=5036790 file=extracted/1008/5036790/2011년 제22회 감정평가사 1차시험 기출문제/2교시.hwp
%% conversion: kordoc 4.24.51 -> tools/gen1.py (auto); review=proofread-sonnet; flags=none
\begin{question}{17}
(주)서울농장은 20×1년 1월 1일에 1년된 돼지 10마리를 보유하고 있다. (주)서울농장은 20×1년 7월 1일에 1.5년된 돼지 5마리를 한 마리당 ￦100,000에 매입하였고, 20×1년 7월 1일에 돼지 6마리가 태어났다. 돼지의 일자별 한 마리당 순공정가치가 다음과 같을 때 (주)서울농장이 동 생물자산과 관련하여 20×1년도 포괄손익계산서의 당기손익에 반영할 평가이익은? (단, 20×1년 중 매각되거나 폐사된 돼지는 없다고 가정함)

{\small\centering
\begin{tabular}{|l|l|r|}
\hline
\multicolumn{1}{|c|}{일 자} & \multicolumn{1}{c|}{내 용} & \makecell{한 마리당\\순공정가치} \\
\hline
20×1. 1. 1 & 1년된 돼지 & ￦80,000 \\
20×1. 7. 1 & 1.5년된 돼지 & 100,000 \\
20×1. 7. 1 & 새로 태어난 돼지 & 50,000 \\
20×1. 12. 31 & 0.5년된 돼지 & 70,000 \\
20×1. 12. 31 & 2년된 돼지 & 130,000 \\
\hline
\end{tabular}\par}

\begin{choices}
  \item ￦300,000
  \item ￦650,000
  \item ￦1,070,000
  \item ￦1,430,000
  \item ￦1,950,000
\end{choices}
\end{question}
\answer{3}
